\section{Conclusion}
\label{sec:conclusion}

Sizing a zonal E/E platform means deciding, together, where the tasks run,
how they are prioritized, how they communicate between cores at LET instants, and how their
messages cross a TSN backbone, under the wiring, memory and safety
constraints of a vehicle and for the whole life of the platform. This
paper stated that problem formally and gave a complete method to solve it.
The timing analyses come with proofs: an optimal priority assignment that
respects precedence, response times with shared-memory contention, bounds
on the maximum reaction time of multi-rate task-chains, and Network
Calculus bounds for credit-based, time-aware and asynchronous shaping and
for frame replication. The heuristic combines bin packing or a new
communication-aware rule with core consolidation, relocation moves, and a
deadline-driven local search with an exact incremental re-analysis. An exact formulation with lazy schedulability cuts shows that
the heuristic is close to optimal, and a lifecycle analysis turns the
provisioning of the hardware into a measurable trade-off with the number
of software updates the platform absorbs.

On applications built from industrial statistics, no bound was exceeded
in a frame-level simulation of the network or in a simulation of the
whole system, and the end-to-end bounds were reached, or nearly; the
packing rules with consolidation came within a few percent of the
optimal allocation in milliseconds on the default platform, where
simulated annealing needs seconds and a genetic algorithm falls behind,
while on harder platforms a genetic algorithm run for tens of seconds,
and the exact formulation, recover most of the gap; communication awareness removed
about two thirds of the inter-ECU communications at the price of a third
more cores; asynchronous traffic shaping removed the burst growth of rings;
strict LET shrank the network bounds but cost deadlines, one producer
period per crossing;
memory contention could add half the cores of a packed mapping; and each
core bought at the start of production bought about half a software
update. Consolidation and the deadline-driven local search account for
most of the benefit of the method, the conclusions hold
when the assumptions of the generator change, and the same method maps
domain, hybrid and distributed architectures, within 5\,\% of the
optimum on the domain and hybrid ones. Measuring that gap also improved
the method: it showed where free software should go on a platform
without a central computer. A whole vehicle of 2000
task-chains is mapped in about a second, so a system architect can compare
architectures, network options and provisioning policies with a
worst-case analysis behind every timing result it reports.

The natural next steps are richer contention models (caches,
interconnects), mixed-criticality scheduling, re-allocation of deployed
software during updates, and an exact treatment of the end-to-end
objective, which this paper improves but does not yet optimize.
